\begin{abstract}
Automated rice leaf pathology classification using deep convolutional neural networks has achieved strong benchmark performance. However, standard classification models predict categorical disease labels without explicit spatial evidence identifying specific lesion regions, frequently learning spurious background and border correlations. In this work, we present RiceGuard, an empirical study evaluating whether weak spatial supervision using bounding-box lesion annotations can provide spatial lesion localization capability while assessing its impact on disease classification and explainability. Built upon an EfficientNet-B0 backbone, RiceGuard jointly optimizes six-class disease classification and $7\times 7$ grid-based lesion localization. A controlled refinement procedure (Phase 3B) addresses foreground-background objectness imbalance via capped positive class weighting ($\text{pos\_weight}=10.00$) and validation-only confidence calibration ($\tau=0.60$, Top-$K=3$). Evaluated on the RiceLeafDiseaseBD benchmark, the Phase 2B classification baseline achieves 90.12\% accuracy and 0.8891 Macro F1, while the Phase 3B multi-task model achieves 88.96\% accuracy, 0.8751 Macro F1, a Localization F1 of 0.0905, and a Mean Matched IoU of 0.6423. Calibration reduces the healthy leaf false positive rate from 21.52\% to 10.97\% and decreases border attention from 18.5\% to 9.2\%. Quantitative Grad-CAM validation against independent RiceSeg5932 segmentation masks demonstrates a slight directional gain in Pointing Game accuracy on mutually correctly classified samples (38.21\% vs. 36.79\%), while all-sample continuous alignment metrics reflect cross-domain distribution shifts. The findings demonstrate that bounding-box multi-task learning successfully imparts spatial localization capability and suppresses background shortcut reliance, accompanied by a measurable, limited classification trade-off.
\end{abstract}
